% !TeX root = ../../infdesc.tex
\section{Sets}
\secbegin{secSets}

\begin{exercise}
Think about why the characterisations of `$\forall x \in X$' in terms of `$\forall x$' and of `$\exists x \in X$' in terms of `$\exists x$' involve different logical operators ($\Rightarrow$ for $\forall$, and $\wedge$ for $\exists$), and verify de Morgan's laws for quantifiers (\Cref{thmDeMorganQuantifiers}) for the expanded forms.
\end{exercise}

\begin{solutions*}
When proving a proposition of the form $\forall x \in X,\, p(x)$, we let $x \in X$ and then prove $p(x)$---thus, we introduce a variable $x$, we assume $x \in X$, and we prove $p(x)$. This is equivalent to introducing an element $x \in \mathcal{U}$ and proving $x \in X \Rightarrow p(x)$, which explains the implication operator in $\forall x,\, x \in X \Rightarrow p(x)$.

When proving a proposition of the form $\exists x \in X,\, p(x)$, we define a specific object $a$, we prove $a \in X$, and we prove $p(a)$. This is equivalent to defining $a \in \mathcal{U}$ and then proving $a \in X \wedge p(a)$, which explains the conjunction operator in $\exists x,\, x \in X \wedge p(x)$.

We can use maximal negation (as in \Cref{secLogicalEquivalence}) to verify de Morgan's laws:
\begin{align*}
\neg \forall x \in X,\, p(x) &\equiv \neg \forall x,\, (x \in X \Rightarrow p(x)) && \\
&\equiv \exists x,\, \neg (x \in X \Rightarrow p(x)) && \text{by \Cref{thmDeMorganQuantifiers}} \\
&\equiv \exists x,\, (x \in X \wedge (\neg p(x))) && \text{by \Cref{exNegationOfImplication}} \\
&\equiv \exists x \in X,\, \neg p(x)
\end{align*}
and

\begin{align*}
\neg \exists x \in X,\, p(x) &\equiv \neg \exists x,\, (x \in X \wedge p(x)) && \\
&\equiv \forall x,\, \neg (x \in X \wedge p(x)) && \text{by \Cref{thmDeMorganQuantifiers}} \\
&\equiv \forall x,\, (x \in X \Rightarrow (\neg p(x))) && \text{by \Cref{exNegationOfConjunctionAlternative}} \\
&\equiv \forall x \in X,\, \neg p(x)
\end{align*}
\end{solutions*}

\begin{exercise}[Russell's paradox has been resolved!]
Repeat \Cref{exRussellsParadox} with the updated definitions of sets and set-builder notation.
\end{exercise}

\begin{solutions*}
Note that by the new characterisation of set-builder notation, we have $R \in R$ if and only if $R \in \mathcal{U} \wedge R \not\in R$.

If $R \in \mathcal{U}$, this means that $R \in R$ if and only if $R \not\in R$, which is a contradiction for the same reason as in \Cref{exRussellsParadox}.

It follows that $R \not\in \mathcal{U}$; but then all this means is that $R \in R$ is false.
\end{solutions*}


\begin{exercise}
Let $a,b,c,d \in \mathbb{R}$ with $a<b$ and $c<d$. Prove that $[a,b) \subseteq (c,d]$ if and only if $a > c$ and $b \le d$.
\end{exercise}

\begin{solutions*}
$(\Rightarrow)$ Assume $[a,b) \subseteq (c,d]$.
\begin{itemize}
\item Since $a \in [a,b)$, we have $a \in (c,d]$ and so $a > c$ as required.
\item Towards a contradiction, assume $b > d$. Then $\frac{b+d}{2} < \frac{b+b}{2} = b$, so $\frac{b+d}{2} \in [a,b)$, and so $\frac{b+d}{2} \in (c,d]$. This implies that $\frac{b+d}{2} \le d$; however since $b>d$ we have $\frac{b+d}{2} > \frac{d+d}{2} = d$, contradicting $\frac{b+d}{2} \le d$. So indeed $b \le d$, as required.
\end{itemize}
$(\Leftarrow)$ Assume $a>c$ and $b \le d$, and let $x \in [a,b)$, so that $x \in \mathbb{R}$ and $a \le x < b$. We need to prove that $x \in (c,d]$. Well, combining these inequalities, we have
\[ c < a \le x < b \le d \]
and so $c < x \le d$, so that $x \in (c,d]$, as required. \qed
\end{solutions*}


\begin{exercise}
Prove that $\{ x \in \mathbb{R} \mid x^2 < x \} = (0,1)$.
\end{exercise}

\begin{solutions*}
Write $A = \{ x \in \mathbb{R} \mid x^2 < x \}$; we prove that $A = (0,1)$ by double containment.
\begin{itemize}
\item $(\subseteq)$ Let $x \in A$. Then $x \in \mathbb{R}$ and $x^2 < x$, so that $0 < x-x^2 = x(1-x)$. Now, if $x \le 0$ then $1-x > 0$, and so $x(1-x) \le 0$; and if $x \ge 1$ then $x > 0$ and $1-x \le 0$, and so $x(1-x) \le 0$. Thus we must have $0<x<1$, and so $x \in (0,1)$, as required.
\item $(\supseteq)$ Let $x \in (0,1)$. Then $x \in \mathbb{R}$ and $0<x<1$. Since $x>0$, the inequality $x<1$ is preserved when we multiply through by $x$, and so $x^2 \le x$. Hence $x \in A$, as required.
\end{itemize}
By double containment, we have $A = (0,1)$. \qed
\end{solutions*}

\begin{exercise}
Prove by double containment that $\{ 0, 1 \} = \{ 1, 0 \}$ and $\{ 0, 0 \} = \{ 0 \}$.
\end{exercise}

\begin{solutions*}
Proving $\{ 0, 1 \} = \{ 1, 0 \}$:
\begin{itemize}
\item $(\subseteq)$ Let $x \in \{ 0, 1 \}$. Then either $x = 0 \in \{ 1, 0 \}$ or $x = 1 \in \{ 1, 0 \}$; in both cases we see that $x \in \{ 1, 0 \}$.
\item $(\supseteq)$ Identical to $(\subseteq)$, but swap the roles of $0$ and $1$.
\end{itemize}
Hence $\{ 0, 1 \} = \{ 1, 0 \}$ by double containment. \qed

Proving $\{ 0, 0 \} = \{ 0 \}$:
\begin{itemize}
\item $(\subseteq)$ Let $x \in \{ 0, 0 \}$. Then $x = 0$ or $x = 0$\dots{} so $x = 0 \in \{ 0 \}$.
\item $(\supseteq)$ Let $x \in \{ 0 \}$. Then $x = 0 \in \{ 0, 0 \}$.
\end{itemize}
Hence $\{ 0, 0 \} = \{ 0 \}$ by double containment. \qed

These proofs are a little bit silly, but they demonstrate that in list notation for sets, we can change the order that the elements are listed and we can repeat elements in the list without changing the set.
\end{solutions*}


\begin{exercise}
Let $a,b \in \mathbb{R}$. Prove that $[a,b]$ is empty if and only if $a > b$, and that $(a,b)$ is empty if and only if $a \ge b$.
\end{exercise}



\begin{exercise}
Let $E$ be an empty set and let $p(x)$ be a predicate with one free variable $x$ with domain of discourse $E$. Show that the proposition $\forall x \in E,\, p(x)$ is true, and that the proposition $\exists x \in E,\, p(x)$ is false. What does the proposition $\forall x \in E,\, x \ne x$ mean in English? Is it true?
\hintlabel{exEverythingIsTrueOfElementsOfEmptySet}{%
Recall from the beginning of \Cref{secSets} that $\forall x \in X,\, p(x)$ is equivalent to $\forall x,\, (x \in X \Rightarrow p(x))$ and $\exists x \in X,\, p(x)$ is equivalent to $\exists x,\, (x \in X \wedge p(x))$. What can be said about the truth value of $x \in E$ when $E$ is empty?
}
\end{exercise}


\begin{exercise}
\label{exEmptySetElementSubsetOfEverySet}
Is it true or false that $\varnothing \in X$ for all sets $X$? Is it true or false that $\varnothing \subseteq X$ for all sets $X$?
\end{exercise}

\begin{exercise}
Explain why the sets $\varnothing$, $\{ \varnothing \}$ and $\{ \{ \varnothing \} \}$ are different, and then determine the truth values of each of the following propositions:
\begin{multicols}{3}
\begin{enumerate}[(a)]
\item $\varnothing \in \varnothing$
\item $\varnothing \subseteq \varnothing$
\item $\varnothing \in \{ \varnothing \}$
\item $\varnothing \in \{ \{ \varnothing \} \}$
\item $\{ \varnothing \} \subseteq \{ \varnothing \}$
\item $\{ \varnothing \} \subseteq \{ \{ \varnothing \} \}$
\end{enumerate}
\end{multicols}
Find more examples of simple elementhood and subsethood propositions---that is, propositions of the form $A \in B$ or $A \subseteq B$---involving the three sets $\varnothing$, $\{ \varnothing \}$, $\{ \{ \varnothing \} \}$; you should find at least one true and one false elementhood proposition, and at least one true and one false subsethood proposition.
\end{exercise}


\begin{exercise}
Write out the elements of $\mathcal{P}(\{1, 2, 3\})$.
\end{exercise}

\begin{exercise}
Let $X$ be a set. Show that $\varnothing \in \mathcal{P}(X)$ and $X \in \mathcal{P}(X)$.
\end{exercise}

\begin{exercise}
Write out the elements of $\mathcal{P}(\varnothing)$, $\mathcal{P}(\mathcal{P}(\varnothing))$ and $\mathcal{P}(\mathcal{P}(\mathcal{P}(\varnothing)))$.
\end{exercise}


\begin{exercise}
Determine, with proof, whether or not each of the following statements is true.
\begin{enumerate}[(a)]
\item $\mathcal{P}(\varnothing) \in \{ \varnothing \}$;
\item $\mathcal{P}(\varnothing) \in \mathcal{P}(\mathcal{P}(\varnothing))$;
\item $\mathcal{P}(\mathcal{P}(\varnothing)) \in \{ \varnothing, \{ \varnothing, \{ \varnothing \} \} \}$.
\end{enumerate}
Repeat the exercise with all instances of `$\in$' replaced by `$\subseteq$'.
\end{exercise}